\Needspace{8\baselineskip}
\section{Worker I: discover jointly inductive polynomial ideals}
\label{sec:ideals}

\subsection{From conservation to preservation of a zero set}

Suppose a state is $x=(x_1,\ldots,x_n)$ and the transition is a polynomial map
$T:\Q^n\to\Q^n$. A conserved polynomial satisfies the identity $p\circ T=p$.
An invariant equation only needs
\[
 p(x)=0\quad\Longrightarrow\quad p(T(x))=0.
\]
These are substantially different requirements. For $p=x^2-y$ and
$T(x,y)=(x^2,y^2)$,
\begin{equation}
 p(T(x,y))=x^4-y^2=(x^2+y)(x^2-y).
 \label{eq:nonlinear-curve}
\end{equation}
The equation $p=0$ is preserved, but the value of $p$ generally changes. A
constant-multiplier ansatz would still miss this example; its multiplier really
is a polynomial in the state.

A family of equations can be mutually supporting. If $T(x,y)=(y,x)$, then the
family $p_1=x$, $p_2=y$ is preserved on the zero state, although neither generator
is conserved. The appropriate acceptance language is
\begin{equation}
 p_i(T_e(x))=
 \sum_{j=1}^{r}A_{eij}(x)p_j(x)+
 \sum_{\ell=1}^{q_e}B_{ei\ell}(x)g_{e\ell}(x),
 \label{eq:ideal-step}
\end{equation}
where edge $e$ is enabled only when all its equality guards
$g_{e\ell}(x)=0$ hold. The matrices $A_e$ and $B_e$ have polynomial entries.

Forge's existing structural section already identifies this more general
identity and correctly warns that simultaneously discovering the $p_j$ and
$A_{eij}$ gives a bilinear problem~\cite{forge-structural}. The algorithm below
does not solve that simultaneous system by pretending it is linear. It freezes
a candidate space, computes ideals generated by that known space, and then
solves a genuinely linear preimage problem. Repeating this operation removes
unsupported candidates until the space stabilizes.

\subsection{Input fragment and semantics}

The executed prototype accepts a finite family of polynomial initial maps
\[
 I_a:\Q^{m_a}\longrightarrow\Q^n
\]
and finitely many guarded polynomial transitions $(T_e,G_e)$, where
$G_e=\{g_{e1},\ldots,g_{eq_e}\}$. It seeks a proof that a supplied goal polynomial
$h$ vanishes at every state reachable from any $I_a(t)$ by enabled transitions.
Zero-parameter maps represent concrete initial states. Multiple initial maps
represent a finite union of parameterized sets.

The initial-map description is a sufficient, explicit source fragment. It does
not accept an arbitrary algebraic initial variety and silently replace that
variety by sampled points. For a Lean application, the front end must prove
that the source initial states are covered by these maps. A map allowed to
range over all rationals may be an over-approximation of natural initial
parameters; that loses power but not validity for the covered source states.

The prototype treats equality guards only. An inequality guard can be omitted
as an over-approximation, so long as the resulting stronger step obligation is
actually proved. Exploiting inequalities requires a different certificate,
perhaps using Forge's existing ordered-polynomial worker; it is not achieved by
turning an inequality into a polynomial equation.

All search is over $\Q$. The resulting identities can be evaluated in a
commutative $\Q$-algebra, such as $\R$. Transport to natural-number code requires
proved cast and update equations. Natural subtraction is not rational
subtraction, and truncating integer division is not field division.

\subsection{The descending-space algorithm}

Let $P_d\subset\Q[x_1,\ldots,x_n]$ be the vector space of polynomials of total
degree at most $d$. Its dimension is $N=\binom{n+d}{d}$. Start with
\begin{equation}
 V_0=\{p\in P_d:\ p(I_a(t))=0\text{ identically for every }a\}.
 \label{eq:initial-space}
\end{equation}
A monomial basis of $P_d$ makes this a rational kernel computation: substitute
each initial map into each monomial, collect coefficients in its own parameter
ring, stack the resulting linear equations, and take the nullspace. Separate
initial maps remain separate blocks; their coefficients must not be added to
each other.

For a subspace $V$, let $J_e(V)=\ideal{V,G_e}$ be the polynomial ideal generated
by a basis of $V$ and the equality guards of edge $e$. Define
\begin{equation}
 \Phi(V)=\{p\in V:\ p\circ T_e\in J_e(V)\text{ for every }e\}.
 \label{eq:ideal-operator}
\end{equation}
Then compute $V_{k+1}=\Phi(V_k)$ until it stabilizes.

For each edge, compute a Gr\"obner basis of the \emph{known} ideal $J_e(V_k)$
using a fixed global monomial order. Reduction modulo that basis defines a
$\Q$-linear map
\[
 L_{e,k}:V_k\to\Q[x],\qquad
 p\longmapsto\NF_{J_e(V_k)}(p\circ T_e).
\]
Although the codomain polynomial ring is infinite-dimensional, only finitely
many monomials occur in the images of the finite basis. Collecting their
coefficients gives a finite matrix. The joint kernel of these matrices is
$V_{k+1}$.

\par\noindent\begin{minipage}{\linewidth}
\begin{lstlisting}
V := kernel of all initial-substitution coefficient matrices
repeat:
    residuals := []
    for each transition e:
        G := GroebnerBasis(basis(V) ++ equalityGuards(e))
        residuals.append([normalForm(p(T_e), G) for p in basis(V)])
    K := rationalKernel(coefficientsOfEachResidualBlock(residuals))
    W := linearCombinations(basis(V), columns(K))
    if dimension(W) = dimension(V): break with V := W
    V := W
extract step multipliers against basis(V) ++ guards(e)
extract goal multipliers against basis(V)
return a candidate only if all independently checked identities hold
\end{lstlisting}
\end{minipage}\par

The implementation uses dimension equality because $W\subseteq V$ is known by
construction. No equality of pretty-printed basis expressions is required.
A production implementation can put each basis in a canonical row-reduced
coefficient form to improve caching, but canonicalization is not needed for the
termination argument.

\begin{theorem}[Finite descent]
In the absence of resource cutoffs, the iteration has at most $\dim V_0$
strictly decreasing rounds. At termination its output $V_*$ satisfies
$p\circ T_e\in\ideal{V_*,G_e}$ for every $p\in V_*$ and every edge $e$.
\end{theorem}
\begin{proof}
The operator maps subspaces to subspaces because substitution, quotient
projection, and coefficient extraction are linear. Also $\Phi(V)\subseteq V$.
Every strict inclusion of finite-dimensional vector spaces lowers dimension by
at least one. When dimensions agree, the subspaces agree; the defining
membership conditions then hold throughout $V_*$.
\end{proof}

The bound counts \emph{outer rounds}, not arithmetic operations, Gr\"obner
pairs, intermediate monomials, or coefficient bits. An individual round can be
expensive. The prototype therefore has distinct basis, polynomial, pair, and
round limits. Hitting one raises a resource-cutoff result; it does not produce
an empty fixed point or a non-invariance verdict.

\subsection{What kind of completeness is obtained?}

\begin{definition}
An admissible degree-$d$ inductive generator space is a subspace $W\subseteq P_d$
such that $W\subseteq V_0$ and
\[
 p\circ T_e\in\ideal{W,G_e}\qquad(p\in W,\ e\text{ an edge}).
\]
Polynomial multipliers have no degree cap in this mathematical definition.
\end{definition}

\begin{theorem}[Greatest admissible space]
With exact completed ideal computations and no resource cutoff, $V_*$ contains
every admissible degree-$d$ inductive generator space.
\end{theorem}
\begin{proof}
Let $W$ be admissible. Initially $W\subseteq V_0$. Suppose $W\subseteq V_k$.
For $p\in W$ and every $e$,
\[
 p\circ T_e\in\ideal{W,G_e}\subseteq\ideal{V_k,G_e}.
\]
Thus $p\in\Phi(V_k)=V_{k+1}$. Induction gives $W\subseteq V_k$ at every round,
and therefore $W\subseteq V_*$ when the iteration stops.
\end{proof}

This theorem is stronger than ``some template happened to work'', but weaker
than completeness for semantic polynomial invariance. Vanishing on a zero set
need not coincide with membership in its defining ideal; radical consequences
and real-algebraic consequences can require additional reasoning. An invariant
may also need generators of degree greater than $d$, a better state encoding,
or a non-polynomial assertion. The initial maps and equality-guard fragment
impose further restrictions.

The final goal is accepted when $h\in\ideal{V_*}$ and a multiplier representation
is extracted. The degree of $h$ itself need not be at most $d$. A high-degree
goal can follow from low-degree generators. Conversely, a goal outside this
ideal does not become false. The result is simply unknown for this route.

\subsection{Multiplier extraction without trusting the Gr\"obner engine}

A normal-form result of zero is useful for search but is not a suitable
standalone replay artifact. The output must say how the reduced polynomial is
built from the original generators.

The prototype uses an extended Buchberger calculation. For original generators
$f_1,\ldots,f_s$, each current basis element carries a polynomial vector
$u_i\in\Q[x]^s$ satisfying
\[
 b_i=\sum_j u_{ij}f_j.
\]
If an S-polynomial is $a b_i-c b_j$, its accompanying vector is
$a u_i-c u_j$. Subtracting a multiple $q b_\ell$ during reduction subtracts
$q u_\ell$ from the vector. Normalizing a nonzero remainder to monic form
normalizes its vector by the same scalar. These operations retain exact
representations throughout the computation.

After reducing a requested polynomial $f$ to zero,
$f=\sum_i q_i b_i$, substitute the recorded vectors to obtain
$f=\sum_j(\sum_i q_i u_{ij})f_j$. Only the latter vector is serialized. The
replayer expands the resulting polynomial identity using an independently
implemented sparse polynomial representation. It does not rerun Buchberger and
does not trust the claimed basis to be Gr\"obner, reduced, minimal, or even
linearly independent.

For larger inputs, a Sage/Singular backend is a sensible replacement for this
small educational search. Sage documents Gr\"obner bases and ideal membership
through its Singular-backed multivariate polynomial interfaces~\cite{sage-ideals}.
The adapter must also request original-generator representations and validate
their orientation by multiplication. A solver's matrix convention is not part
of Forge's mathematics: the adapter must prove that its decoded multipliers
satisfy exactly Equation~\eqref{eq:ideal-step}.

\subsection{The finite certificate and its soundness}

For invariants $p_1,\ldots,p_r$, the certificate contains their rational
coefficient lists, the multiplier rows $A_{eij},B_{ei\ell}$ for each transition,
and goal multipliers $c_1,\ldots,c_r$. The original problem is supplied
\emph{separately}. Acceptance verifies
\begin{align}
 p_i(I_a(t))&=0 &&\text{for every }i,a,\label{eq:ideal-init-check}\\
 p_i(T_e(x))&=\sum_jA_{eij}p_j+\sum_\ell B_{ei\ell}g_{e\ell}
 &&\text{for every }i,e,\\
 h(x)&=\sum_jc_j(x)p_j(x).\label{eq:ideal-goal-check}
\end{align}
Every equality is a formal polynomial identity, not an equation checked on
selected points.

\begin{theorem}[Invariant-certificate soundness]
If the finite checks above hold, then $h(x)=0$ at every reachable state of the
encoded transition system.
\end{theorem}
\begin{proof}
At an initial state, Equation~\eqref{eq:ideal-init-check} makes every generator
zero. Suppose they are zero at a reachable state and edge $e$ is enabled. Every
term on the right of Equation~\eqref{eq:ideal-step} vanishes, using the current
invariant equations and the edge's guards. Hence every generator is zero after
the transition. Induction on the number of transitions proves this conjunction
at every reachable state. Finally Equation~\eqref{eq:ideal-goal-check} makes the
goal vanish.
\end{proof}

The proof uses a conjunction of all generators as its induction predicate.
There is no circular admission of one conjectured lemma by another: the
\emph{whole conjunction} has an independently checked initial case and step.
This is an important difference from installing mutually dependent auxiliary
lemmas as if each were already known.

\subsection{Examples that discriminate between mechanisms}

\paragraph{Scaling versus conservation.}
Initialize $(x,y)=(t,t^2)$ and use $T=(2x,4y)$. The worker finds $p=x^2-y$ and
$p\circ T=4p$. In the degree-two space, a polynomial conserved under the scaling
can have only its constant term: the nonconstant monomials scale by factors
$2^a4^b\ne1$. Initial vanishing then forces that constant to zero. Thus the
conservation subspace is zero while the inductive space is nonzero.

\paragraph{Nonlinear multipliers.}
With the same initial map and $T=(x^2,y^2)$, the discovered multiplier is
$x^2+y$. This example excludes an implementation that merely generalizes
$p\circ T=p$ to $p\circ T=\lambda p$ with a constant $\lambda$.

\paragraph{Several possible updates.}
The same invariant survives both the squaring update and
$T_1=(x+1,y+2x+1)$, since $p\circ T_1=p$. A single invariant family must pass
every encoded edge. Checking only the transition that happened to produce a
candidate would be an unsound source restriction.

\paragraph{Equality guards.}
Initialize $(x,y)=(0,t)$, permit the transition
$(x,y)\mapsto(x+y,y+1)$ only under $y=0$, and seek $x=0$. Its step is the simple
identity $x+y=1\cdot x+1\cdot y$. The second term is justified by the guard,
not by an invariant about all reachable values of $y$. The mutation suite
removes this guard and verifies rejection of the old certificate.

\paragraph{Finite orbits and a genuine degree boundary.}
Initialize $x=0$ and use $T(x)=1-x$. At degree one, $V_0=\spanQ\{x\}$ but the
preimage test removes it: $1-x\notin\ideal{x}$. The search returns unknown for
$x^2-x=0$. At degree two, the dimensions are $2,1,1$, and the fixed space is
$\spanQ\{x^2-x\}$. This example explains why an empty low-degree result does not
refute the source theorem.

\subsection{Implementation refinements worth making next}

The first production optimization should be a good exact coefficient backend,
not an enlarged template vocabulary. Keep sparse coefficient vectors, cache the
normal forms of basis images within a round, and reuse initial-substitution
matrices across degree increases. A stable monomial order makes these caches
meaningful. Incremental Gr\"obner algorithms may exploit related ideals between
rounds, but dropping generators can invalidate a previously computed basis, so
``reuse'' must not mean retaining invalid reductions.

For a program with several control-flow locations, use a separate bounded
space $V_\ell$ at each location. For an edge $e:\ell\to m$, require every target
generator $p\in V_m$ to satisfy $p\circ T_e\in\ideal{V_\ell,G_e}$. The same
monotone-descent argument works on the finite product of the location spaces;
the total dimension bounds strict decreases. This extension is a design here,
not part of the executed prototype. It is particularly attractive for mutually
recursive functions whose invariants naturally differ by control location.

A second refinement is goal-directed generator selection. Run the complete
bounded-space pass at a modest degree, then prune a successful certificate to
the generators reachable from its goal-multiplier dependencies. Recheck all
remaining step identities after pruning. It is safe to shrink a proved
certificate; it is not safe to remove a generator during search merely because
its name does not occur in the final goal.
